% Pursuing Stacks (cotes 134-2 to 134-8): archive cote 134-4, site page 40 (PDF page 41), typescript page 128.
% typescript. Draft compared: ps3.tex:1998-2061.
% AI-transcribed by Claude Opus 5.5 (claude-opus-5-5, adaptive thinking, effort high), Message Batches API, batch ps2 (msgbatch_01NVwr1Yuo8qJf5rykeGBXPK), 2026-11-12; image review recorded in records/CHECK_DECISIONS_134_3_4.tsv.
\scanpage{134-4/40}
\mspage{128}

structure, and the two are associated. This comes from the fact that
both compositions of the functors (2),(3) must be the identity functors
(being compatible with the ``under M'' structure), hence \hand{\ensuremath{\alpha}} and \hand{\ensuremath{\beta}} are isomor%
phisms, which shows both R \hand{\ensuremath{\in}} Q$_{\rm o}$ and W \hand{\ensuremath{\in}} L$_{\rm o}$ , and that R and L are associa-%
ted. In terms of the set theoretic situation (1), this may be described
by using the order relations on Q and L, and the fact that r and s are
monotone maps, and satisfy moreover
(*) \qquad $sr(q)\hand{\ensuremath{\leq}} q$ \;,\; $rs(\hand{\ensuremath{\ell}})\hand{\ensuremath{\leq}}\hand{\ensuremath{\ell}}$ \quad (any $q\hand{\ensuremath{\in}}Q$, $\hand{\ensuremath{\ell}}\hand{\ensuremath{\in}}L$) \;,
which implies that the set C$_{\rm o}$ of pairs $(q\unsure{_{,}}\,, \hand{\ensuremath{\ell}})$ of associated elements
of Q$_{\rm o}$,L$_{\rm o}$ can be described also as
\[ C_{o} = \hand{\ensuremath{\{}}(q,\hand{\ensuremath{\ell}})\hand{\ensuremath{\in}} Q\text{x}L \hand{\ensuremath{\mid}} \hand{\ensuremath{\ell}}\hand{\ensuremath{\leq}}r(q),\quad q\hand{\ensuremath{\leq}} s(\hand{\ensuremath{\ell}})) \qquad . \]

Thus it does'nt seem evident to get a homotopy structure on M, just
starting with an\deleted{y} R \hand{\ensuremath{\in}} Q or \xout{any} a W \hand{\ensuremath{\in}} L, without assuming beforehand that
R is a homotopy relation, or W a homotopism structure. T$\overset{\hand{\scriptscriptstyle\perp\perp}}{\underset{\hand{-}}{\mathrm h}}$e condition
$C_{1,2}$ on page 115 may be viewed as a condition on a pair (R,W)\hand{\ensuremath{\in}} QxL,
and it clearly implies
\[ R \hand{\ensuremath{\leq}} s(W) \quad ; \]
if we assume moreover \quad W=r(R) \quad we get \quad R\hand{\ensuremath{\leq}} s(r(R)) \quad and hence, in view
of the first inequality (*) above ,
\[ R = sr(R) \quad \text{i.e.} \quad R \hand{\ensuremath{\in}} Q_{o} \quad , \]
i.e. R is a homotopy relation, as asserted somewhat quickly yesterday.

T$\overset{\hand{\scriptscriptstyle\perp\perp}}{\mathrm h}$e only standart way for getting homotopy structures in a general
category M which I \xout{\illegible{1}} can see by now , is in terms of an arbitrary
set \hand{\ensuremath{\Sigma}}$_{h}$ of intervals in M (assuming only that M \over{a}{a}dmits finite products).
As soon as the focus gets upon intervals for describing homotopy structu-%
res, the situation becomes typically non-autodual - in contrast to
Quillen's autodual treatment of homotopy relations (via``left'' and ``right''
homotopies, involving respectively ``cylinder'' and ``path'' - objects). \xout{Thi}
This is in keeping with the \inserted{\unsure{highly|highty} non auto-dual} axiom\unsure{\over{s}{s}} on universal disjoint \uline{sums in M,} which
we finally introduced by the end of yesterday's \uline{reflection.}\diagram{\begin{tikzpicture}[baseline=-.3ex,x=1em,y=1ex]\draw[color=blue,line width=.5pt] (0,0)--(.12,2.5);\draw[color=blue,line width=.5pt] (.13,0)--(.25,2.5);\end{tikzpicture}} To come
back however upon the relationships between the four basic ``homotopy
notions'' introduced in yesterday's notes, I would now rather symbolize
these relations in the following diagram of maps between sets

\deleted{\illegible{1}}

\margin{40}
